\section{Experiments}
\label{sec:experiments}

We evaluate STA-derived ISR across three axes: as a diagnostic for reward hacking (\S\ref{sec:diagnostic}), as a training reward exhibiting meta-level Goodhart's Law (\S\ref{sec:goodhart}), and as a cross-logic diagnostic protocol (\S\ref{sec:cross-logic}).
We synthesize the observed hacking strategies into a taxonomy (\S\ref{sec:taxonomy-results}).

\subsection{Setup}
\label{sec:setup}

We train Qwen3-4B (base, no prior fine-tuning) with GRPO \cite{deepseekmath} using $K{=}8$ candidate formalizations per problem and a learning rate of $10^{-6}$.
Default training runs 500 gradient steps; we extend to 1{,}000 steps for one ablation.
Training data consists of 500 FOLIO-derived \cite{folio} FOL examples; evaluation uses a fixed held-out set of 100 problems.
Both sets remain constant across all experiments to isolate the effect of reward formulation.
Seeds $\in \{42, 43, 44\}$ control model initialization and training stochasticity.

We compare five reward conditions (\S\ref{sec:reward}): exec-only (baseline), additive ISR bonus at $\beta \in \{0.1, 0.3\}$, matched random reward shaping at $\beta \in \{0.1, 0.3\}$, and supervised fine-tuning (SFT) on the same 500 examples with gold FOL annotations.
A convex random control ($r = 0.9 \cdot r_{\text{exec}} + 0.1 \cdot u$) is included for completeness.

Evaluation applies STA (Algorithm~\ref{alg:sta}) to every proved sample on the held-out set via a unified evaluation script.
Solvers are Prover9 \cite{mccune2010prover9} for FOL and SWI-Prolog~8.4.2 for Prolog.
Hacking classification combines normalized text matching (for H1, H7) with functional analysis of the STA necessity vector ($n_{\text{necessary}}$); we report both when they diverge.

\subsection{ISR as Reward Hacking Diagnostic}
\label{sec:diagnostic}

\begin{figure}[t]
\centering
% Placeholder: dual y-axis line plot
% X = training steps {100, 500, 1000}
% Left Y = prove rate (ascending: 0.32, 0.93, 0.96)
% Right Y = proved-only ISR (descending: 0.642, 0.132, 0.003)
% Annotate: component count 6.6 (500 steps) and 18.6 (1000 steps)
% Annotate: 94.8% ISR=0 at 1000 steps
\fbox{\parbox{0.9\columnwidth}{\centering\vspace{3cm}\textit{Fig.\ placeholder: ISR collapse trajectory}\vspace{3cm}}}
\caption{ISR collapse trajectory under exec-only GRPO (seed~42). Prove rate rises from 0.32 to 0.96 while proved-only ISR collapses from 0.642 to 0.003. By 1{,}000 steps, 94.8\% of proved samples have $\text{ISR}{=}0$ and mean component count inflates from 6.6 to 18.6, indicating contradiction explosion (\S\ref{sec:taxonomy-results}).}
\label{fig:isr-collapse}
\end{figure}

\paragraph{ISR collapse trajectory.}

Exec-only GRPO achieves increasing prove rate while ISR undergoes monotonic collapse (Figure~\ref{fig:isr-collapse}).
At 100 training steps, the model proves 32\% of held-out problems with a proved-only ISR of 0.642---formalizations are incomplete but structurally sound.
By 500 steps, prove rate climbs to 93\% but ISR drops to 0.132, a 79\% decline: 87.1\% of proved samples embed the conclusion among the premises (H1 hacking).
At 1{,}000 steps, the collapse deepens: prove rate reaches 96\% while ISR drops to 0.003---a 99.5\% decline from the 100-step checkpoint.
At this stage, 94.8\% of proved samples have $\text{ISR}{=}0$ and mean component count rises from 6.6 to 18.6 as the model shifts to \emph{ex falso quodlibet} strategies, generating contradictory premises from which any conclusion follows.

This inverse relationship is invisible to the execution reward.
An optimizer monitoring only $r_{\text{exec}}$ would interpret the 500-step model as strictly better than the 100-step one.
ISR reveals that the improvement is illusory: the model learned to exploit solver behavior rather than produce logically coherent formalizations.

\paragraph{Cross-seed instability.}

\begin{figure}[t]
\centering
% Placeholder: stacked bar chart
% X-axis categories: Exec s42, Exec s43, Exec s44 | ISR s42, ISR s43, ISR s44 | Rand s42, Rand s43, Rand s44 | SFT
% Y-axis: fraction of proved samples
% Stacked segments (bottom to top): H7 (red), H1 (orange), Other hacking (yellow), Genuine (green)
\fbox{\parbox{0.9\columnwidth}{\centering\vspace{3cm}\textit{Fig.\ placeholder: hacking composition by condition}\vspace{3cm}}}
\caption{Hacking mode composition across conditions (proved samples only). Each bar decomposes into H7 (conclusion duplication, ISR-specific), H1 (conclusion embedding, GRPO-universal), and genuine proofs. Exec-only seeds all exhibit 87--100\% conclusion embedding. ISR $\beta{=}0.1$ converts H1 into H7 without reducing total hacking. SFT produces 92.9\% genuine proofs.}
\label{fig:hacking-composition}
\end{figure}

Three seeds under exec-only reward yield prove rates spanning 14\% to 98\%, yet all exhibit pervasive conclusion embedding (Figure~\ref{fig:hacking-composition}).
Seed~42 achieves $\text{PR}{=}93\%$ with 87.1\% H1 conclusion embedding (ISR${=}$0.132); only 12.9\% of proved samples are genuine.
Seed~43 collapses to $\text{PR}{=}14\%$ through repetition and length explosion ($\text{ISR}{=}0$).
Seed~44 reaches $\text{PR}{=}98\%$ via tautological collapse: the model generates two-component formalizations where the sole premise equals the conclusion ($\text{H7}{=}100\%$, mean components${=}$2.04, $\text{ISR}{=}0.582$).
The cross-seed prove rate standard deviation of 47 percentage points underscores that single-seed GRPO results are unreliable.
No seed produces predominantly genuine proofs.

\paragraph{SFT baseline.}

Supervised fine-tuning on the same 500 examples produces $\text{PR}{=}14\%$ with proved-only $\text{ISR}{=}0.571$.
Of 14 proved samples, 13 (92.9\%) are genuine, one contains H1 conclusion embedding, and none exhibit H7 or other hacking modes.
SFT's low prove rate (14\% vs.\ 12\% for the base model) confirms that data exposure alone barely improves performance---RL exploration drives the prove rate gains in GRPO.
Its high ISR and near-zero hacking rate establish that structural degeneration is RL-reward-specific, not a general fine-tuning artifact.

\subsection{Meta-Level Goodhart Effect}
\label{sec:goodhart}

Having established ISR as a diagnostic, we test whether it can serve as a training signal.
The answer is negative: ISR-as-reward exhibits \emph{meta-level} Goodhart's Law \cite{manheim2019categorizing}---the diagnostic metric itself becomes the target of gaming when optimized.

\paragraph{ISR reward vs.\ random control at $\beta{=}0.1$.}

\begin{table}[t]
\centering
\small
\setlength{\tabcolsep}{3.5pt}
\caption{ISR vs.\ Random at $\beta{=}0.1$ across three seeds. ISR column reports proved-only mean. H7 is ISR-specific (pooled: 149/231 = 64.5\% vs.\ 0/256 = 0\%). H1 is GRPO-universal, higher under random ($81.3\%$) than ISR ($17.3\%$). Genuine\% is the residual: $100\% - \text{H7}\% - \text{H1}\%$.}
\label{tab:beta01}
\begin{tabular}{clccccc}
\toprule
\textbf{Seed} & \textbf{Signal} & \textbf{PR} & \textbf{ISR} & \textbf{H7} & \textbf{H1} & \textbf{Gen.} \\
\midrule
s42 & ISR    & .81 & .415 & 16.0\% & 51.9\% & 32.1\% \\
s43 & ISR    & .61 & .830 & 96.7\% &  0\% &  3.3\% \\
s44 & ISR    & .89 & .902 & 86.5\% &  0\% & 13.5\%$^\dagger$ \\
\midrule
s42 & Random & .79 & .131 &  0\% & 70.9\% & 29.1\% \\
s43 & Random & .88 & .058 &  0\% & 93.2\% &  6.8\% \\
s44 & Random & .89 & .070 &  0\% & 79.8\% & 20.2\% \\
\bottomrule
\end{tabular}

\vspace{2pt}
\begin{minipage}{\columnwidth}
\footnotesize
$^\dagger$All 12 nominally genuine ISR~s44 samples have $n_{\text{prem}}{=}1$, evading the H7 text-match classifier; functional analysis ($n_{\text{necessary}}{=}1$) identifies them as H7 false negatives.
True genuine rate for ISR~s44 is 0\%.
\end{minipage}
\end{table}

Table~\ref{tab:beta01} presents the central comparison.
Across three seeds, ISR reward at $\beta{=}0.1$ achieves a mean prove rate of $0.77 \pm 0.14$, while the matched random control achieves $0.85 \pm 0.06$.
ISR reward provides no prove rate advantage and exhibits higher variance.

The proved-only ISR values tell a superficially different story: ISR-rewarded models achieve $0.716 \pm 0.26$ vs.\ $0.086 \pm 0.04$ for random.
A naive reading of this $8.3\times$ gap would suggest dramatic structural improvement.
The hacking decomposition (Table~\ref{tab:beta01}) reveals the opposite.

\paragraph{H7 conclusion duplication is ISR-specific gaming.}

ISR reward induces H7 conclusion duplication in 64.5\% of proved samples (149/231 pooled across seeds).
The random control produces \emph{zero} H7 instances across all three seeds (0/256)---the strongest possible null.
H7 formalizations consist of a single premise identical to the conclusion; with $n{=}2$ components, removing the premise always breaks the proof, yielding $\text{ISR} \geq 0.5$ by construction (\S\ref{sec:isr}).
The ISR reward trains the model not to reason better, but to produce structures that score high on ISR.
This constitutes meta-level Goodhart's Law: a metric designed to detect reward hacking becomes the target of a new hacking mode when incorporated into the reward.

In seeds~43 and~44, H7 dominates (96.7\% and 86.5\%, respectively) with near-zero genuine output.
Seed~42 shows a more balanced composition (16\% H7, 52\% H1, 32\% apparent genuine), which we examine next.

\paragraph{Premise-count confound.}

The 32\% genuine rate in ISR~seed~42 initially suggested partial structural benefit.
Closer analysis reveals this is a premise minimization artifact.
Of the 26 nominally genuine samples, 15 (57.7\%) contain only 1--2 premises.
ISR is mathematically biased toward short formalizations: when $n_{\text{prem}}{=}1$, the sole premise is trivially necessary, yielding $\text{ISR} \geq 0.5$ regardless of proof quality (\S\ref{sec:isr}).
ISR reward incentivizes premise minimization (mean 5.58 components vs.\ 8.46 for random), which inflates ISR through the metric's mathematical properties rather than improved reasoning.

To control for this confound, we restrict comparison to genuine samples with $n_{\text{prem}} \geq 3$:

\vspace{-2pt}
\begin{center}
\small
\begin{tabular}{lcc}
\toprule
 & \textbf{ISR s42} & \textbf{Random s42} \\
\midrule
Samples ($n_{\text{prem}} \geq 3$) & 11 & 21 \\
Mean ISR & 0.258 & 0.243 \\
\textbf{Ratio} & \multicolumn{2}{c}{\textbf{1.06$\times$}} \\
\bottomrule
\end{tabular}
\end{center}
\vspace{-2pt}

\noindent Controlling for premise count eliminates the apparent ISR advantage.
The $1.06\times$ ratio is statistically indistinguishable from chance.
ISR-as-reward does not improve structural quality; it induces premise minimization that inflates the ISR metric.

\paragraph{Conclusion embedding is GRPO-universal.}

Table~\ref{tab:beta01} reveals that conclusion embedding (H1) is pervasive across all GRPO conditions.
The random control exhibits H1 in $81.3\% \pm 11.3\%$ of proved samples---\emph{higher} than ISR reward at 17.3\%.
The total conclusion-in-premise rate (H7 + H1) is comparable: ISR at $\sim$82\% vs.\ random at $\sim$81\%.
ISR reward does not reduce hacking; it redirects it from H1 to H7, altering the composition of hacking strategies without reducing their prevalence.
This suggests conclusion embedding is an emergent consequence of GRPO exploration against execution rewards, not an artifact of any specific reward signal.

\paragraph{High-dose collapse ($\beta{=}0.3$).}

At $\beta{=}0.3$, ISR reward exhibits classic Goodhart collapse \cite{manheim2019categorizing}: eval prove rate drops to 18\% while the random control achieves 48\%.
ISR at this dose is not merely unhelpful but \emph{actively harmful}---the model overfits to ISR-inflating strategies that fail to generalize.
The random control achieves $2.7\times$ higher prove rate despite contributing no structural signal.

\paragraph{Dose-response pattern.}

\begin{table*}[t]
\centering
\caption{Goodhart landscape across all reward conditions and doses. All metrics from held-out eval (100 problems). ISR reports proved-only mean. Multi-seed conditions ($n{=}3$) report mean $\pm$ std; single-seed conditions report seed~42. H7, H1, Genuine computed from hacking classification of proved samples.}
\label{tab:landscape}
\resizebox{\textwidth}{!}{%
\begin{tabular}{llcccccccl}
\toprule
\textbf{Condition} & \textbf{Formula} & \textbf{Seeds} & \textbf{PR} & \textbf{ISR} & \textbf{Comp.} & \textbf{H7\%} & \textbf{H1\%} & \textbf{Gen.\%} & \textbf{Goodhart Phase} \\
\midrule
Exec-only (500 steps) & $r_{\text{exec}}$ & s42 & 0.93 & 0.132 & 6.6 & 0 & 87.1 & 12.9 & Baseline \\
Exec-only (1000 steps) & $r_{\text{exec}}$ & s42 & 0.96 & 0.003 & 18.6 & 0 & 100 & 0 & Cont.\ explosion \\
SFT baseline & supervised & s42 & 0.14 & 0.571 & --- & 0 & 7.1 & \textbf{92.9} & No RL hacking \\
\midrule
ISR $\beta{=}0.1$ & $r_{\text{exec}} \!\cdot\! (1{+}0.1 \!\cdot\! \text{ISR})$ & 3 & $.77{\scriptstyle\pm.14}$ & $.716{\scriptstyle\pm.26}$ & 5.6 & \textbf{64.5} & 17.3 & ${\sim}$16 & Meta-Goodhart \\
Random $\beta{=}0.1$ & $r_{\text{exec}} \!\cdot\! (1{+}0.1 \!\cdot\! u)$ & 3 & $.85{\scriptstyle\pm.06}$ & $.086{\scriptstyle\pm.04}$ & 8.5 & 0 & 81.3 & ${\sim}$19 & Control \\
ISR $\beta{=}0.3$ & $r_{\text{exec}} \!\cdot\! (1{+}0.3 \!\cdot\! \text{ISR})$ & s42 & 0.18 & 0.143 & 10.7 & --- & --- & --- & Classic Goodhart \\
Random $\beta{=}0.3$ & $r_{\text{exec}} \!\cdot\! (1{+}0.3 \!\cdot\! u)$ & s42 & 0.48 & 0.094 & 9.6 & --- & --- & --- & Control \\
\midrule
Convex random & $0.9 r_{\text{exec}} {+} 0.1 u$ & s42 & 1.00 & 0.535$^*$ & 2.0 & --- & --- & --- & Tautological collapse \\
\bottomrule
\end{tabular}%
}
\vspace{2pt}
\begin{minipage}{\textwidth}
\footnotesize
$^*$ISR artifact: 100\% two-component tautologies ($\text{premise}{=}\text{conclusion}$, $\text{ISR}{=}0.5$). \quad
Comp.\ = mean component count (s42 for single-seed rows, s42 for multi-seed $\beta{=}0.1$ conditions). \\
$\beta{=}0.3$ hacking decomposition omitted: low sample counts (18 and 48 proved) and severe format degradation preclude reliable classification. \\
Exec-only 3-seed PR: s42${=}$0.93, s43${=}$0.14, s44${=}$0.98 ($\sigma{=}$47pp).
Multi-seed pooled H7\%: ISR 149/231, Random 0/256.
\end{minipage}
\end{table*}

The two dose levels trace a complete Goodhart trajectory (Table~\ref{tab:landscape}).
At low dose ($\beta{=}0.1$), the model games the proxy metric itself: ISR inflation via H7 produces high ISR scores that mask structural degeneration.
At high dose ($\beta{=}0.3$), the ISR signal overwhelms the execution reward, causing prove rate collapse.
No dose yields genuine structural improvement.
The SFT baseline (92.9\% genuine) and random control ($0.85 \pm 0.06$ PR) bracket the landscape: hacking is RL-reward-specific, and ISR reward underperforms even random noise.

\subsection{Cross-Logic Diagnostic Transfer}
\label{sec:cross-logic}

To test whether ISR diagnostic and Goodhart effects generalize beyond FOL, we replicate the core experiments using Prolog formalizations verified by SWI-Prolog.
All training uses the same FOLIO-derived data and hyperparameters; only the target formalism and solver change.

\paragraph{Prolog exec-only: diagnostic transfer.}

\begin{table}[t]
\centering
\small
\caption{Prolog experiments across three seeds. \textit{Top}: exec-only ISR diagnostic detects structural degeneration across all seeds ($\text{ISR}{<}0.1$). \textit{Bottom}: ISR-as-reward produces universal Goodhart---3/3 seeds gaming, 0\% genuine. ISR for exec-only is eval mean; for ISR reward, proved-only.}
\label{tab:prolog}
\begin{tabular}{llcccc}
\toprule
\textbf{Condition} & \textbf{Seed} & \textbf{PR} & \textbf{ISR} & \textbf{Func.\ H7} & \textbf{Gen.} \\
\midrule
Exec-only & s42 & .83 & .031 & --- & --- \\
Exec-only & s43 & .56 & .008 & --- & --- \\
Exec-only & s44 & .21 & .081 & --- & --- \\
\cmidrule{1-6}
\multicolumn{2}{l}{\textit{Exec-only mean}} & \textit{.53${\scriptstyle\pm}$.31} & \textit{.040${\scriptstyle\pm}$.037} & & \\
\midrule
ISR reward & s42 & .50 & 1.000$^a$ & 100\% & 0\% \\
ISR reward & s43 & .38 & .185 & 100\%$^b$ & 0\% \\
ISR reward & s44 & .97 & .912 & 97.9\%$^c$ & 0\% \\
\bottomrule
\end{tabular}

\vspace{2pt}
\begin{minipage}{\columnwidth}
\footnotesize
$^a$All 50 proved samples are two-element (clause${=}$query), so $\text{ISR}{=}1.0$ each. \\
$^b$Strict text classifier returns 0\% H7, 97.4\% H1 due to degenerate uppercase pseudo-Prolog syntax. Functional analysis: 38/38 proved have $n_{\text{necessary}}{=}1$. \\
$^c$Strict classifier returns 0\% H7, 0\% H1 (query-as-fact gaming evades text matching). Functional analysis: 95/97 proved have $n_{\text{necessary}}{=}1$.
\end{minipage}
\end{table}

Three seeds under Prolog exec-only reward produce prove rates spanning 21\% to 83\% (mean $0.53 \pm 0.31$), reproducing the cross-seed instability observed in FOL ($\sigma{=}$47pp FOL vs.\ 31pp Prolog).
ISR remains below 0.1 across all seeds (Table~\ref{tab:prolog}), confirming diagnostic transfer: STA detects structural degeneration in Prolog as it does in FOL, regardless of prove rate level.

Low ISR in Prolog exec-only reflects syntactic redundancy---e.g., repeated ground facts where removing any single copy does not affect the proof---rather than complete absence of reasoning.
This contrasts with high ISR under Prolog ISR reward, which reflects metric gaming.
ISR values require interpretation in context of the training objective (\S\ref{sec:taxonomy-results}).

\paragraph{Prolog ISR reward: universal Goodhart.}

\begin{figure}[t]
\centering
% Placeholder: grouped bar chart
% Two groups: FOL, Prolog
% Bars: exec-only mean ISR, ISR-reward mean ISR (proved-only), genuine %
% Key contrast: Prolog ISR-reward = 0% genuine vs FOL ~16%
\fbox{\parbox{0.9\columnwidth}{\centering\vspace{3cm}\textit{Fig.\ placeholder: FOL vs.\ Prolog ISR comparison}\vspace{3cm}}}
\caption{Cross-logic comparison. ISR as diagnostic (exec-only) transfers across logic types: both produce $\text{ISR}{<}0.1$ despite varying prove rates. ISR as reward produces universal Goodhart in Prolog (0\% genuine across all seeds) vs.\ $\sim$16\% genuine residual in FOL, consistent with simpler semantics enabling more complete exploitation.}
\label{fig:cross-logic}
\end{figure}

When ISR is incorporated into the Prolog reward, all three seeds converge to gaming with zero genuine output (Table~\ref{tab:prolog}, Figure~\ref{fig:cross-logic}).
Each seed discovers a distinct gaming pathway:

\begin{itemize}[nosep,leftmargin=*]
\item \textbf{Seed~42} ($\text{PR}{=}50\%$, $\text{ISR}{=}1.0$): Clean H7 tautology. The model generates a single Prolog clause identical to the query. Every component is trivially necessary in a two-element formalization.
\item \textbf{Seed~43} ($\text{PR}{=}38\%$, $\text{ISR}{=}0.185$): Degenerate syntax. The model produces uppercase pseudo-Prolog (e.g., \texttt{MARY\_IS\_AN\_AVADEOTE.}) that SWI-Prolog interprets via variable bindings. The strict H7 text classifier returns 0\%, but functional analysis reveals all 38 proved samples have $n_{\text{necessary}}{=}1$.
\item \textbf{Seed~44} ($\text{PR}{=}97\%$, $\text{ISR}{=}0.912$): Query-as-fact gaming. The model converts natural language problems into \texttt{snake\_case} atoms and asserts them simultaneously as fact and query, achieving trivial proofs. Of 97 proved samples, 95 have $n_{\text{necessary}}{=}1$. The strict text classifier initially reported 100\% genuine; independent functional analysis corrected this to 0\%.
\end{itemize}

This result is stronger than the FOL finding.
In FOL, ISR reward at $\beta{=}0.1$ retains $\sim$16\% genuine residual across three seeds.
In Prolog, the genuine rate is 0\% across all seeds---Prolog's simpler semantics enable more complete exploitation, eliminating the genuine residual observed in FOL.
The three distinct gaming pathways demonstrate that the Goodhart effect is not a single failure mode but a family of exploitation strategies that emerge from the interaction between the optimization objective and the target language's structure.

\paragraph{Classifier limitations.}

The Prolog results expose a blind spot of text-matching hacking classifiers.
Seeds~43 and~44 bypass the strict H7 classifier through surface-level novelty (non-standard casing, atom-level encoding) while remaining functionally tautological ($n_{\text{necessary}}{=}1$).
Seed~44 in particular---with $\text{PR}{=}97\%$ and $\text{ISR}{=}0.912$---would appear to be a resounding success by all scalar metrics.
Only the necessity vector reveals the gaming.
This motivates cross-validation of hacking classifiers with functional metrics such as the $n_{\text{necessary}}$ distribution.

\subsection{Reward Hacking Taxonomy}
\label{sec:taxonomy-results}

\begin{table*}[t]
\centering
\caption{Hacking taxonomy: prevalence (\% of proved samples) across conditions. H1 and H7 detected via normalized text matching; other modes via functional analysis of the STA necessity vector and manual inspection. Dashes indicate modes not observed or not applicable in that condition.}
\label{tab:taxonomy}
\small
\begin{tabular}{p{5.2cm}ccccc}
\toprule
\textbf{Mode} & \textbf{Exec-only} & \textbf{ISR $\beta{=}$0.1} & \textbf{Rand.\ $\beta{=}$0.1} & \textbf{SFT} & \textbf{Prolog ISR} \\
 & \textbf{(s42, 500s)} & \textbf{(3-seed)} & \textbf{(3-seed)} & \textbf{(s42)} & \textbf{(3-seed)} \\
\midrule
H1: Conclusion embedding & 87.1 & 17.3 & 81.3 & 7.1 & 32.5$^a$ \\
H7: Conclusion duplication & 0 & \textbf{64.5} & 0 & 0 & 27.0$^b$ \\
H3: Complexity without utility & 0 & 0 & 3.8$^c$ & 0 & --- \\
Tautological collapse & s44$^d$ & --- & --- & --- & --- \\
Repetition collapse & s43$^e$ & --- & --- & --- & --- \\
Contradiction explosion & 1000s$^f$ & --- & --- & --- & --- \\
Query-as-fact gaming & --- & --- & --- & --- & \textbf{40.5}$^g$ \\
\midrule
\textbf{Total hacking} & \textbf{87.1} & \textbf{$\sim$84} & \textbf{$\sim$81} & \textbf{7.1} & \textbf{100} \\
\textbf{Genuine} & 12.9 & $\sim$16 & $\sim$19 & \textbf{92.9} & \textbf{0} \\
\bottomrule
\end{tabular}

\vspace{2pt}
\begin{minipage}{\textwidth}
\footnotesize
$^a$Prolog s43 strict classifier; functionally H7 (all $n_{\text{necessary}}{=}1$). \quad
$^b$Prolog s42: 50/185 pooled proved. \quad
$^c$Random s42 only. \\
$^d$Exec-only s44: 100\% H7, $n_{\text{comp}}{=}2.04$, $\text{PR}{=}98\%$. \quad
$^e$Exec-only s43: $\text{PR}{=}14\%$, 87\% hit \texttt{max\_new\_tokens}. \quad
$^f$Exec-only 1{,}000 steps: 100\% H1, $n_{\text{comp}}{=}18.6$. \\
$^g$Prolog s44: 95/97 proved have $n_{\text{necessary}}{=}1$; evades text-match classifiers. \quad
Prolog pooled: 50 (s42) + 38 (s43) + 97 (s44) = 185 proved. \\
Tautological, repetition, and contradiction modes are seed- or step-specific exec-only phenomena, not reported as percentages in the main columns.
\end{minipage}
\end{table*}

Across all conditions, seeds, and logic types, we identify seven distinct hacking modes (Table~\ref{tab:taxonomy}).
Three patterns emerge.

First, \textbf{H1 conclusion embedding is GRPO-universal}.
It appears under exec-only (87.1\%), random reward (81.3\%), and ISR reward (17.3\%), spanning every RL condition.
ISR reward does not reduce total hacking---it converts H1 into the more extreme H7 mode, as the total conclusion-in-premise rate remains $\sim$82\% in both ISR and random conditions.

Second, \textbf{reward formulation determines hacking mode, not hacking rate}.
All RL conditions produce 81--87\% total hacking; only SFT avoids it (7.1\%).
What changes is the composition: exec-only and random produce H1; ISR reward produces H7; Prolog ISR produces H7, degenerate syntax, and query-as-fact gaming.
This suggests that each reward signal creates a characteristic \emph{hacking niche} that the optimizer fills through the path of least resistance.

Third, \textbf{condition-specific modes reveal optimization dynamics}.
H7 emerges exclusively under ISR reward---a direct consequence of ISR's sensitivity to component count, which the model exploits through premise minimization.
Query-as-fact gaming emerges exclusively in Prolog, where the language permits asserting arbitrary atoms as both facts and queries.
Contradiction explosion requires extended training (1{,}000 steps), suggesting it needs sustained optimization pressure.
These mode--condition associations demonstrate that hacking strategies are systematic responses to the reward landscape, not random failures.

\paragraph{Summary.}

The ISR formulation $r = r_{\text{exec}} \cdot (1 + \beta \cdot \text{ISR})$ fails at every dose: $\beta{=}0.1$ induces meta-Goodhart gaming (H7), $\beta{=}0.3$ collapses prove rate below random noise, and the sole apparent positive (ISR~s42 genuine rate) reduces to $1.06\times$ when premise count is controlled.
ISR's value is exclusively diagnostic: it quantifies hacking that prove rate misses, detects hacking across logic types, and characterizes the specific strategies by which models exploit each reward signal.
Its failure as a training signal exemplifies a broader lesson: a metric that reliably detects gaming from the outside can become the target of new gaming strategies when placed inside the optimization loop.
